% ============================================================================
% REMERCIEMENTS
% ============================================================================
\chapter*{Remerciements}
\addcontentsline{toc}{chapter}{Remerciements}

Ce mémoire de fin d'études, couronnement de deux années de formation intensive en Master Sciences et Techniques, n'aurait pu aboutir sans le concours et le soutien de nombreuses personnes que je tiens à remercier sincèrement.

Je tiens en premier lieu à exprimer ma profonde gratitude envers mon encadrant académique, \textbf{Pr. Kaoutar ALLABOUCHE}, pour son encadrement rigoureux, ses conseils avisés et sa disponibilité constante. Ses remarques pertinentes et son expertise scientifique ont été déterminantes pour l'orientation et la qualité de ce travail.

Mes remerciements s'adressent également à mon encadrant industriel, \textbf{M. Mohamed EDDOUCH}, Responsable Maintenance chez SEWS Maroc, pour sa confiance, son accompagnement professionnel et le partage de son expertise technique. Je remercie également \textbf{M. Hamza BAHLOULI} pour son support technique précieux et sa patience lors des phases d'implémentation.

Je souhaite remercier l'ensemble du personnel de \textbf{SEWS Maroc}, particulièrement l'équipe du département Maintenance et de l'Atelier Test Électrique, pour leur accueil chaleureux, leur collaboration active et leur contribution à la réussite de ce projet.

Mes remerciements vont également aux membres du jury qui ont accepté d'évaluer ce travail et de l'enrichir par leurs remarques constructives.

Je remercie l'administration de la \textbf{Faculté des Sciences de l'Université Ibn Tofail} et l'ensemble du corps professoral du département Génie Informatique pour la qualité de la formation dispensée et les compétences académiques et professionnelles acquises durant ce cursus.

Enfin, je ne saurais oublier de remercier ma famille pour leur soutien indéfectible, leurs encouragements permanents et leur patience tout au long de mon parcours universitaire.

\vspace{1cm}
\begin{flushright}
\textit{Aissam MALKOUT}\\
\textit{Kénitra, juillet 2026}
\end{flushright}
\clearpage
